\documentclass[11pt]{amsart}
\usepackage{pjm}

\author{John Teague}
\title{Surface Finder}

\begin{document}

\maketitle

\begin{abstract}
	We present a practical algorithm to enumerate every PL-embedded surface of a particular type in a given 4-manifold triangulation. We are primarily interested in finding closed surfaces and surfaces whose boundaries are contained within the 3-manifold boundary of the triangulation. In the latter case, if this 3-manifold triangulation is a triangulation of $S^3$, we present an algorithm to computationally identify the PL-embedded links associated to these boundaries. We also include experimental results which (hopefully) verify the slice genus of every knot under 10 crossings.
\end{abstract}

\section{Introduction}

Recall that, in dimension 4, the smooth and PL categories are equivalent by work of Cairns and Whitehead \ref{}; results which hold for PL 4-manifolds also hold for smooth 4-manifolds. Thus, one can study the smooth topology of 4-manifolds from an entirely PL perspective. This more combinatorial setting lends itself to more tractable approaches by computational techniques and experimental verification. 

Similarly, one can approach the study of smoothly embedded surfaces in 4-manifolds from this point of view. This subsumes topics like link concordance and other related notions in knot theory. To show that two links are concordant, it suffices to demonstrate a PL surface between PL copies of those links. So, one could show that a given knot is slice by computationally exhibiting a disk which the knot bounds in $B^4$.

In practice, these approaches take place in \textit{triangulations} of 4-manifolds. Although one typically uses ``triangulation'' to mean ``simplicial complex'', we use the term here to refer to \textit{generalized triangulations}, also known as unordered $\Delta$-complexes. These are more computationally flexible than simplicial complexes since, for example, faces of the same simplex can be glued together.



\section{Preliminaries}

Recall definition of sliceness, slice genus, etc.

\section{Surface Finding Algorithm}

\subsection{Knot complement algorithm}

\section{Computational Results}

\section{Results}

\newpage

% \bibliographystyle{alpha}
% \bibliography{refs}

\end{document}
